\chapter{Lista 3.}	
\begin{problem}[1]

	Niech $G\times X\rightarrow G$ będzie dowolną funkcją. Definiujemy $\psi: G \rightarrow X^{X}$ przez $\psi(g)(x) := g\cdot x$.

	Dowieść, że następujące warunki są równoważne.
	\begin{enumerate}[label=\alph*)]
		\item $\cdot $ jest lewym działaniem G na X.
		\item $\psi[G] \subseteq S_{X}$ oraz $\psi: G\rightarrow S_{X}$ jest homomorfizmem.
	\end{enumerate}
\end{problem}
\begin{solution}

	TODO		
\end{solution}
\begin{problem}[2]

	Niech DA(G,H). będzie zbiorem wszystkich działań G na H przez automorfizmy, a Hom(G, Aut(H)) grupą wszystkich homomorfizmów z G w grupę Aut(H). Dowieść, że funkcje
	 \[
	F: \text{DA(G,H)}\rightarrow \text{Hom(G, Aut(H)) oraz } F': \text{Hom(G, Aut(H))} \rightarrow \text{DA(G,H)}
	.\] 
	zadane wzorami $F(\cdot )(g)(h) = g\cdot h$ oraz $F'(\psi)(g,h) = \psi(g)(h)$ są wzajemnie odwrotnymi bijekcjami.

	\textbf{Wsk.} Wywnioskować, że wystarczy udowodnić wzajemne zawierania dziedziny i przeciwdziedziny  oraz udowodnić te zawierania.
\end{problem}
\begin{solution}

	TODO	
\end{solution}
\begin{problem}[3]

	Dwa elementy $a,b \in G$ nazywamy sprzężonymi, gdy istnieje $g \in G$, takie żę $gag^{-1}=b$. Zauważyć, że elementy a i b są sprzężone wtedy i tylko wtedy, kiedy leżą w tej samej orbicie względem działania grupy na sobie przez automorfizmy wewnętrzne. Wywnioskować, że relacja sprzężenia jest relacją równoważności, a jej klasy abstrakcja są orbitami tego działania: 

	klasa abstrakcji a to orbita a, czyli zbiór $\{gag^{-1}: g\in G\} = \{g^{-1}ag\} =: a^{G}$, czyli jest to klasa sprzężoności a(zgodnie z terminologią z wykładu).
\end{problem}
\begin{solution}

	TODO		
\end{solution}
\begin{problem}[4]

	\begin{enumerate}[label=\alph*)]
		\item Znaleźć i udowodnić kryterium na to, kiedy dwa elementy w grupie $S_{n}$ są sprzężone. Jako wniosek opisać klasy sprzężoności w grupie $S_{n}$.
		\item Które z poniższych permutacji są sprzężone? W przypadku gdy dwie permutacje są sprzężone, znaleźć element, który je sprzęga.
			\[
			       \begin{pmatrix} 1 & 2 & 3 & 4 & 5 & 6 \\
				 2 & 4 & 5 & 1 & 3 & 6 
				\end{pmatrix} \text{ oraz } 
					\begin{pmatrix} 
			         1 & 2 & 3 & 4 & 5 & 6 \\
		                 4 & 5 & 1 & 3 & 6 & 2\end{pmatrix} 
			.\] 

			\[
				\begin{pmatrix} 1 & 2 & 3 & 4 & 5 & 6 &7 \\
				5 & 7 & 2 & 1 & 6 & 4 & 3 
			\end{pmatrix} 
			\text{ oraz }
					\begin{pmatrix} 1 & 2 & 3 & 4 & 5 & 6 & 7 \\
					6 & 7 & 5 & 3 & 4 & 2 & 1
				\end{pmatrix} 
			.\] 
	\end{enumerate}
\end{problem}
\begin{solution}

	TODO
\end{solution}
\begin{solution}
	
	TODO
\end{solution} 
\begin{problem}[5]
	
	Udowodnić, że wszystkie automorfizmy  $S_3$ są wewnętrzne. 
\end{problem} 
\begin{solution}
	
	TODO
\end{solution} 
\begin{problem}[6]
	
	Niech G będzie grupą. Załóżmy, że istnieje $g\in G$, taki że rząd(g)$\neq 1,2$. Udowodnić, że Aut(G)$\neq \{\id_{G}\} $
\end{problem} 
\begin{solution}
	
	TODO
\end{solution} 
\begin{problem}[7]
	
	Załóżmy, ze grupa G działa na zbiorze X. Dowieść, że stabilizatory elementów leżących w jednej orbicie mają równą moc.

	\textbf{Wsk.} Dla elementóœ x,y z tej samej orbity wyrazić $G_{y}$ w terminach $G_{x}$.
\end{problem} 
\begin{solution}
	
	TODO
\end{solution} 
\begin{problem}[8]
	
	Opisać orbity działania $GL_{n\left( \R \right) }$ na $\R^{n}$.
\end{problem} 
\begin{solution}
	
	TODO
\end{solution} 
\begin{problem}[9]
	
	Niech $\left( A,+ \right) $ będzie grupą przemienną. Udowodnić, że poniższy wzór
	\[
	\forall_{a\in A}\quad  0\cdot a = a, 1\cdot a = -a 
	.\] 

	zadają działanie $\Z_{2}$ na A przez automorfizmy. Wskazać odpowiadający temu działaniu $/Psi: \Z_2 \rightarrow \text{AutAa)}$. Kiedy $\Psi$ jest monomorfizmem?
\end{problem} 
\begin{solution}
	
	TODO
\end{solution} 
\begin{problem}[10]
	
	Wyznaczyć centrum $S_3$ i centrum $D_4$.
\end{problem} 
\begin{solution}
	
	TODO
\end{solution} 
\begin{problem}[11]
	
	Niech $H\le G.$ Udowodnić bezpośrednio (tzn. nieodwołując się do tw. Lagrange'a), że $|G /H| = |H \setminus G|$. 
\end{problem} 
\begin{solution}
	
	TODO
\end{solution} 

